\section*{Đánh giá trước khi chỉnh sửa}

Mã nguồn báo cáo đã đáp ứng phần lớn yêu cầu CO5173, nhưng bản hiện tại
chưa sẵn sàng để nộp. Đánh giá dưới đây tập trung vào cấu trúc, nội dung,
minh chứng nghiệp vụ và mức độ hoàn thiện của các yêu cầu BR.

\begin{table}[htbp]
\centering
\small
\begin{tabularx}{\textwidth}{@{}>{\bfseries}p{0.19\textwidth}Y@{}}
\toprule
Tiêu chí & Đánh giá \\
\midrule
Cấu trúc &
Mã LaTeX có sáu mục chính, phù hợp với rubric. Tuy nhiên, PDF hiện có vẫn
dùng bố cục cũ (\code{main.tex} và \code{main.pdf}); cần biên dịch lại sau
khi chỉnh sửa. \\
Nội dung kỹ thuật &
Báo cáo đã trình bày bối cảnh, người dùng, nguồn dữ liệu, kiến trúc, công
nghệ, ứng dụng và kết luận. Báo cáo cũng nêu phương án thay thế như Flink
cho Spark Structured Streaming và so sánh PostgreSQL SQL với Python
(Mục~\ref{sec:solution} và Mục~\ref{chap:eval}). \\
Đánh giá ứng dụng &
Tính đúng đắn và hiệu năng có minh chứng. Hiệu quả sử dụng được minh họa
qua việc giảng viên khoanh vùng ca kiểm thử lỗi và cố vấn theo dõi cảnh báo.
Tuy vậy, bảng kiểm chứng nghiệp vụ mới liệt kê kịch bản cho BR4--BR8; cần
nêu rõ kết quả demo hoặc giới hạn của BR1--BR3 (Mục~\ref{chap:eval}). \\
Số lượng BR &
Đề yêu cầu tối thiểu hai BR cho mỗi thành viên. Báo cáo có tám BR, đủ nếu
nhóm có bốn người; nếu có năm người thì cần thêm hai BR. Bìa hiện có năm
dòng thông tin thành viên nên chưa xác nhận được quy mô nhóm. Cần đối chiếu
đề bài CO5173 và thông tin trong \code{config/info.tex}. \\
\bottomrule
\end{tabularx}
\end{table}

\subsection*{Các BR cần làm rõ}

Qua đối chiếu thiết kế với mã nguồn, BR2 là khoảng thiếu lớn nhất. BR1 và
BR3 đã có các dịch vụ xử lý Kafka riêng, nhưng chưa chứng minh được đường
tích hợp hoàn chỉnh với giao diện. BR4--BR8 đã có các chức năng chính; riêng
BR7 cần tách rõ thao tác tua lùi khỏi số liệu seek tổng hợp. Các luồng đầu
cuối chưa được chạy lại trong lần rà soát này.

\begin{description}[leftmargin=2.6cm,labelwidth=2.4cm,style=nextline]
  \item[\br{1} -- Thông báo]
  Đã có Kafka consumer, bộ luật và Redis trong Compose riêng. Tuy nhiên,
  API thông báo của giao diện vẫn đọc kho demo, nên còn thiếu đường nối từ
  sự kiện Kafka đến thông báo hiển thị cho người học. Push và email hiện
  mới ghi log.

  \item[\br{2} -- Gợi ý học tập]
  Đã có gợi ý theo luật và bảng đặc trưng dài hạn. Còn thiếu đặc trưng phiên
  theo luồng và việc kết hợp hai nhánh qua feature store để phục vụ gợi ý.
  Đây là khoảng thiếu chính, kể cả khi tiếp tục dùng luật để xếp hạng.

  \item[\br{3} -- Chấm bài]
  Đã có API nhận bài, hàng đợi Kafka, worker và bộ chuẩn hóa kết quả. Code
  Lab trong classroom vẫn gọi bộ chấm demo rồi lưu kết quả vào PostgreSQL.
  Cần bổ sung minh chứng cho luồng nộp bài từ Code Lab qua Kafka và worker,
  rồi trả phản hồi về Code Lab.

  \item[\br{4}--\br{6}]
  Đã có minh chứng cho các chức năng chính: phát hiện ca kiểm thử bất
  thường, cảnh báo từ CDC qua cửa sổ trượt và báo cáo tiến độ có ngưỡng
  nhóm tối thiểu. Báo cáo còn thiếu benchmark BR4/BR5 mới và đánh giá ngoại
  tuyến BR5 trên OULAD.

  \item[\br{7} -- Heatmap]
  Đã có heatmap theo giây và phiên bản video. Đề bài yêu cầu rewind, trong
  khi mart hiện đếm seek chung cho cả tua lùi và tua tới; vì vậy cần bổ sung
  số đếm tua lùi riêng.

  \item[\br{8} -- Nội dung và tương tác]
  Đã có cả hai phần theo thiết kế: đồng bộ cache/chỉ mục tìm kiếm và phát
  hiện xu hướng tương tác giảm.
\end{description}

\subsection*{Ưu tiên trước khi nộp}

Cần sửa phần báo cáo mô tả adapter BR1/BR3 như công việc tương lai: các dịch
vụ tương ứng đã có, còn khoảng trống nằm ở tích hợp với giao diện và minh
chứng đầu cuối. Về số lượng yêu cầu, tám BR đủ cho nhóm bốn người; nhóm năm
người cần bổ sung hai BR theo quy định.

Trước khi nộp, cần chốt số thành viên, điền thông tin bìa, biên dịch lại và
đặt tên PDF theo mã nhóm. Repo đã có benchmark Spark--Flink và đánh giá
ngoại tuyến BR5 trên OULAD, có thể dùng để củng cố Mục~5. Kết quả OULAD cần được trình bày riêng với
demo giả lập. Đồng thời, cần đối chiếu lại các yêu cầu
về nội dung và quy cách nộp trong đề bài CO5173.